% Part: lambda-calculus
% Chapter: lambda-definability
% Section: partial-recursive-functions

\documentclass[../../../include/open-logic-section]{subfiles}

\begin{document}

\olfileid[id]{lam}{ldf}{par}
\olsection{Fungsi Rekursif Parsial yang \usetoken{S}{lambda definable}}

Fungsi rekursif parsial ialah fungsi yang diperoleh dari fungsi-fungsi dasar
melalui komposisi, rekursi primitif, dan minimisasi tak berbatas.
Fungsi-fungsi ini berbeda dari fungsi rekursif umum karena fungsi-fungsi
yang digunakan dalam pencarian tak berbatas tidak disyaratkan bersifat
reguler. Tanpa syarat
keteraturan, fungsi yang didefinisikan melalui minimisasi mungkin tidak
terdefinisi pada sebagian masukan.

Sekilas, tampaknya metode yang sama dengan yang digunakan untuk menunjukkan
bahwa semua fungsi rekursif umum (yang total) bersifat !!{lambda definable}
dapat digunakan untuk membuktikan bahwa semua fungsi rekursif parsial bersifat
!!{lambda definable}. Sebagai contoh, komposisi $f$ dengan $g$
!!{lambda defined} oleh $\lambd[x][F (G x)]$ jika $f$ dan $g$ masing-masing
!!{lambda defined} oleh suku $F$ dan~$G$. Namun, hal ini menimbulkan masalah
ketika fungsi-fungsi tersebut parsial. Jika $g(x)$ tidak terdefinisi, maka
$G x$ tidak mempunyai bentuk normal. Dalam kebanyakan kasus, ini berarti
bahwa $F (G x)$ juga tidak mempunyai bentuk normal, dan itulah yang kita
inginkan. Namun, tinjau kasus ketika $F$ adalah
$\lambd[x][\lambd[y][y]]$; dalam kasus ini, $F (G x)$ justru mempunyai bentuk
normal, yaitu $\lambd[y][y]$.

Masalah ini bukannya tidak dapat diatasi. Ada cara-cara untuk
!!{lambda define} semua fungsi rekursif parsial sedemikian sehingga nilai yang
tidak terdefinisi direpresentasikan oleh suku tanpa bentuk normal. Akan tetapi,
cara-cara tersebut agak lebih rumit dan kurang intuitif daripada pendekatan
yang kita gunakan bagi fungsi rekursif umum. Di sini kita mencatat teoremanya
tanpa bukti:

\begin{thm}
  Semua fungsi rekursif parsial bersifat !!{lambda definable}.
\end{thm}

\end{document}
